Autism spectrum disorder (ASD) is fundamentally characterized by persistent deficits in communication and social interaction. Eye-tracking (ET) technology functions as an objective tool that provides quantitative data on visual attention processes, which serve as the foundation for many cognitive and communication skills. This study introduces a novel framework designed to advance personalized ASD interventions through data-driven innovation. The research involved three school-aged children diagnosed with ASD. We developed an integrated procedure wherein eye-tracking metrics, such as gaze path, fixation time, and area of interest, were utilized to analyze the specific attentional characteristics of each child. This analysis provided objective evidence that enabled professionals to iteratively adjust traditional interventions, specifically the picture exchange communication system (PECS) and the assessment of basic language and learning skills - revised (ABLLS-R), to better align with the cognitive profile and preferences of each participant. Post-intervention assessments indicated substantial improvements in targeted communication behaviors across all three subjects. Furthermore, data analysis established a strong association between the metric-based adjustment of interventions and improved behavioral outcomes. These results demonstrate a robust impact and highlight the value of this approach for educators seeking to enhance therapeutic effectiveness. This study confirms that an eye-tracking-based framework represents a feasible and effective method for quantifying and personalizing interventions for children with autism. The validity of the framework is supported by consistent improvements across participants and corroborated by behavioral data, thereby reinforcing its potential as a reliable instrument for evidence-based decision-making.
